\documentclass[11pt]{article}
\newcommand{\runningtitle}{The optimal torsion window}
\usepackage{lmodern}
\pdfmapfile{+lm.map}
\usepackage[a4paper,top=24mm,bottom=25mm,left=28mm,right=28mm]{geometry}
\usepackage{amsmath,amssymb,booktabs,array,mathtools,enumitem,xcolor,fancyhdr,hyperref}
\definecolor{ink}{HTML}{17364B}
\definecolor{teal}{HTML}{00685E}
\hypersetup{colorlinks=true,linkcolor=teal,urlcolor=teal,citecolor=teal,pdfauthor={OpenAI}}
\urlstyle{same}
\setlength{\parskip}{3.1pt plus 1pt}
\setlength{\parindent}{0pt}
\setlist{itemsep=1pt,topsep=3pt}
\setlength{\headheight}{15pt}
\pagestyle{fancy}\fancyhf{}
\fancyhead[L]{\small\sffamily\color{ink}\runningtitle}
\fancyhead[R]{\small\sffamily\color{ink}Research note}
\fancyfoot[R]{\small\thepage}
\fancyfoot[L]{\small\sffamily\color{ink}Transfers from the Dirichlet GRH proof}
\renewcommand{\headrulewidth}{0.3pt}
\newcommand{\headline}[3]{%
 {\small\sffamily\bfseries\color{teal}#1\par}
 \vspace{6pt}{\raggedright\LARGE\bfseries\color{ink}#2\par}
 \vspace{7pt}{\small\sffamily #3\par}\vspace{8pt}
 {\color{teal}\hrule height 0.9pt}\vspace{8pt}}
\newcommand{\grhlink}[2]{\href{https://github.com/lsernau/grh-lean/blob/72fead102f663d5c7a407964a13dc0cf9d977207/grh/#1}{#2}}
\newcommand{\incominglink}[2]{\href{https://github.com/VladimirReshetnikov/ProveIt/blob/274841c1cb137aafbb60ab5f519c51fd42eee98e/docs/incoming/#1}{#2}}
\setcounter{tocdepth}{1}
\setlength{\emergencystretch}{2em}
\raggedbottom

\begin{document}
\headline{LOCAL ARITHMETIC}{The optimal torsion window}{A sharp specialization theorem for cyclotomic and formal-group torsion \hfill 25 November 2026}
An auxiliary step in the Lean proof chooses a finite-order twist while
controlling the divisibility of a polynomial value. A previous analysis
obtained zero loss at a fixed positive Gauss radius and asked for the
best possible bound on the order. We determine the least upper exponent,
construct an extremal polynomial, and show that one whole coset of
twists always works. The threshold remains exact for bounded integral
power series. The method also determines the answer in every finite
formal height over an unramified base, including supersingular elliptic
curves. Its proof uses local interpolation alone.

\paragraph{An example to fix the scale.}
For $p=3$, allow roots of order at least $9$. A budget of one means that
the least valuation of a coefficient, after the degree penalty is
added, is at most one. At every comparison radius $\rho\ge1/6$ the old
sufficient order bound was $3^{47}$. The optimal order bound is
$3^5=243$. There is a degree-26 polynomial that defeats every permissible
root through order $81$. No polynomial of smaller degree does so, at
any budget.

\paragraph{The source of the question.}
The five conclusions in \grhlink{src/Statements.lean}{Statements} and
\grhlink{src/Main.lean}{Main} concern zeta and Dirichlet characters.
The readable \grhlink{GUIDE.md}{guide} traces the half-plane
estimate back to three comparisons. Short integral theta lifts give
an upper bound on a lattice covolume, a second moment supplies enough
independent forms to use that bound, and twist-dependent lower bounds
control the same determinant with its index correction. In the
\grhlink{src/New/Reduction/MainInequality.lean}{main inequality} the
common leading height cancels. Local twists help control the denominators
of projections in the lower bound. The pertinent
\grhlink{src/New/Congruences/Selmer/RootsOfUnity.lean}{roots-of-unity
lemmas} use a positive tolerance for a finite list of polynomial
evaluations. This note concerns exact control of one valuation, or of
the sum for a product. Zero error for each member of an arbitrary list
is a stronger requirement.

The direct predecessor is the \incominglink{grh_exact_torsion_transfer_2026-11-25.zip}
{exact torsion report}, which explicitly lists optimal constants among
its remaining questions. Other incoming work already gives sharp
unramified and packet results in related formulations. The improvement
here concerns one fixed-radius objective with unrestricted coefficients.
The notes and scripts in this archive are independent arguments and
checks; the upstream \grhlink{README.md}{replay records} concern the
original Lean development. We inspected that source without rebuilding it.

\section{The least possible order for exact torsion specialization}
\label{sec:optimal-torsion}
A twisting step in the Lean proof seeks characters of controlled prime-power order at which polynomial valuations stay small. The incoming report \incominglink{grh_exact_torsion_transfer_2026-11-25.zip}{on exact torsion selection} proved existence at zero tolerance relative to a fixed Gauss valuation. Its uniform exact specialization theorem gave the sufficient upper exponent
\[
  h_0+\left\lceil p^{h_0}B(h_0+B+2)\right\rceil .
\]
The choice of coefficient field affects the best possible exponent. We determine it exactly for unrestricted integral coefficients. For comparison we give the unramified optimum from the earlier reports.

Let $p$ be prime, $v(p)=1$ on $\mathbb C_p$, $v(0)=+\infty$, and $\mathcal O=\mathcal O_{\mathbb C_p}$. Write $U_h$ for the primitive $p^h$-th roots of unity, $e_h=(p-1)p^{h-1}$ and $r_h=1/e_h$. Thus $v(\zeta-1)=r_h$ for $\zeta\in U_h$. Fix $h_0\ge1$, $\rho\ge r_{h_0}$ and $B\ge0$. For nonzero $g(T)=\sum a_jT^j\in\mathcal O[T]$ use
\[
       G_u(g)=\min_j\{v(a_j)+ju\},\qquad A=G_\rho(g).
\]
A root $\zeta$ is good for $g$ if $v(g(\zeta-1))\le A$. The radius $\rho$ is fixed while the order varies. Let $H_{\rm opt}$ be the smallest $H\ge h_0$ guaranteeing a good root on some shell $h_0\le h\le H$ for every $g$ with $A\le B$.

\subsection{An exact threshold, with a full coset of good roots}
\noindent\textbf{Theorem.}
For unrestricted coefficients in $\mathcal O$,
\[
                   \boxed{H_{\rm opt}=h_0+\lfloor p^{h_0-1}B\rfloor.}
\]
More precisely, put $m=p^{h_0-1}$ and $L=H-h_0+1$. If $A-G_0(g)<L/m$, an entire coset of $\mu_m$ consists of good roots of the same order in the prescribed window. The threshold is sharp even for monic polynomials with coefficients in the integers of a finite cyclotomic extension.

\medskip\noindent\emph{Proof.}
Write $n=p^L$, so $mn=p^H$. Fix a primitive $mn$-th root $\xi$ and a primitive $n$-th root $\eta$. The admissible exponents $b$ modulo $mn$ are exactly those with $b\not\equiv0\pmod n$. For each $1\le a<n$, choose any $b_a\equiv a\pmod n$, and set $x_a=\xi^{b_a}-1$. For distinct residues $a,c$,
\[
 v(x_a)=m^{-1}v(\eta^a-1),\qquad
 v(x_a-x_c)=m^{-1}v(\eta^a-\eta^c).
\]
Indeed $b_a$ and $b_a-b_c$ have the same $p$-adic orders as $a$ and $a-c$, respectively, all below $L$. Passing from an exponent modulo $p^L$ to its lift modulo $p^H$ multiplies the order of the corresponding root of unity by $m$.

Let $P(T)=\prod_{a=1}^{n-1}(T-x_a)$. Each $v(x_a)\le r_{h_0}\le\rho$, so Gauss multiplicativity and $\prod_{a=1}^{n-1}(1-\eta^a)=n$ give
\[
          G_\rho(P)=\sum_a v(x_a)=L/m=:D .
\]
Differentiating $(X^n-1)/(X-1)$ at $X=\eta^a$ gives
\[
          v(P'(x_a))=m^{-1}\bigl(L-v(\eta^a-1)\bigr)
                       =D-v(x_a).
\]
It follows that $\ell_a(T)=P(T)/((T-x_a)P'(x_a))$ has $G_\rho(\ell_a)=0$ for every $a$. In particular, interpolation of any polynomial of degree below $n-1$ from these nodes has no valuation loss on the disc of radius $p^{-\rho}$.

By homogeneity, scale $g$ by the reciprocal of a coefficient of minimal valuation so that $G_0(g)=0$ and its coefficients remain integral. The normalized $A$ is then below $D$. Write $g=P Q+R$ with $\deg R<n-1$. Monic division over $\mathcal O$ leaves $Q,R$ integral. Hence $G_\rho(PQ)\ge D>A$ (using $G_u(0)=+\infty$) and $G_\rho(R)=A$. Interpolation gives
\[
               R(T)=\sum_a g(x_a)\ell_a(T),\qquad
               A\ge\min_a v(g(x_a)).
\]
This holds for every choice of the representatives. If each of the $n-1$ residue classes had a bad representative, we could select one from each, contradicting this inequality. At least one class therefore consists entirely of good roots. It is a coset of $\mu_m$ and contains $m$ roots of a common order.

For sharpness take $b_a=a$ and $g=P$. It has $A=D$, $G_0(P)=0$. At a selected node its value is zero. At another admissible root $\xi^b$ let $b\equiv a\pmod n$. Then $v(\xi^b-\xi^a)> r_{h_0}$, since their quotient is a nontrivial $m$-th root of unity. If $m=1$ there are no such extra points. The remaining factors have the same valuations as at $x_a$, since $v(x_a-x_c)\le r_{h_0}$ for $c\ne a$. Thus
\[
        v(P(\xi^b-1))=D-v(x_a)+v(\xi^b-\xi^a)>D .
\]
All roots in the window are bad. Consequently a window of $L$ shells works for every $A\le B$ exactly when $L>mB$. The asserted formula follows, including its integer endpoints.\hfill$\square$

\medskip
For instance, $p=3$, $h_0=2$, $B=1$ gives the optimal upper exponent $5$ in place of the earlier sufficient value $47$. At endpoint $H=4$, the witness polynomial uses $3^3-1=26$ nodes at a primitive $81$st root and has budget exactly $1$. Extending to $H=5$ guarantees three good roots forming a coset. The comparison concerns the order \emph{exponent}; the actual order is $p^h$.

\subsection{Comparison with the known unramified coefficient case}
The incoming manuscripts \incominglink{grh-torsion-and-zero-data.zip}{on torsion and zero data} and \incominglink{grh-torsion-and-prime-sampling.zip}{on torsion and prime sampling} obtained the exceptional-shell count for bounded ramification, the latter giving the precise per-shell divisibility cost below. They treat a stronger simultaneous formulation. We record the scalar specialization and its short proof to clarify the role of the coefficient field.
\noindent\textbf{Theorem.}
For coefficients in the ring of integers of $\mathbb Q_p$ or of any finite unramified extension $K/\mathbb Q_p$,
\[
                        \boxed{H_{\rm opt}=h_0+\lfloor B\rfloor.}
\]
Every primitive root on the selected shell is good. More generally, if $K$ has absolute ramification index $e$, failure on every shell in the window implies
\[
 A-G_0(g)\ \ge\ \sum_{h=h_0}^{H}\frac{1}{\gcd(e,e_h)} .
\]

\medskip\noindent\emph{Proof.}
If a shell has even one bad point, at least two different coefficient terms attain $G_{r_h}(g)$. Otherwise its unique minimizing term dominates every evaluation on that shell and gives $v(g(\zeta-1))=G_{r_h}(g)\le A$. Let $j_-<j_+$ be the extreme minimizing indices. Their equality gives $(j_+-j_-)/e_h\in e^{-1}\mathbb Z$, hence $j_+-j_-$ is divisible by $e_h/\gcd(e,e_h)$. By the Newton polygon theorem the polynomial has exactly $j_+-j_-$ roots of valuation $r_h$, counted with multiplicity. Equivalently, this number is the drop in slope of the concave function $u\mapsto G_u(g)$ at $r_h$, as follows by factoring $g=b\prod_\alpha(T-\alpha)$ and using
\[
                 G_u(g)=v(b)+\sum_\alpha\min(u,v(\alpha)).
\]
Each such root contributes exactly $r_h$ to $G_\rho(g)-G_0(g)$, and different shells use disjoint roots. Every root contributes nonnegatively. Summation proves the necessary inequality. In fact if the right side exceeds the budget, at least one shell has no bad point at all. For $e=1$ this costs one unit per shell. Conversely the polynomial $\prod_{h=h_0}^H\Phi_{p^h}(1+T)$ lies in $\mathbb Z[T]$, vanishes on the entire window, and has $G_\rho=L$, by its roots. This proves optimality.\hfill$\square$

\subsection{Power series, several variables, and determinant packets}
For $s\ge1$ and $F=\sum_{\nu}a_\nu\mathbf T^\nu\in\mathcal O[[T_1,\ldots,T_s]]$ use $G_\rho(F)=\min_\nu(v(a_\nu)+\rho|\nu|)$. The minimum exists for nonzero $F$. With $G_\rho(F)\le B$, the unrestricted optimal upper exponent above applies to each coordinate. There exist $\zeta_j$ of orders in the window with $v(F(\boldsymbol\zeta-\mathbf1))\le G_\rho(F)$. To see this first treat a polynomial. Expand in the last $s-1$ variables and choose a coefficient $f_\beta(T_1)$ with $G_\rho(f_\beta)+\rho|\beta|=G_\rho(F)$. Specialize $T_1$ at a good root for $f_\beta$. The remaining polynomial has integral coefficients and Gauss valuation no greater than the former value. Continue one variable at a time. For a series, truncate to total degree $d$ with $(d+1)\min(\rho,r_H)>B$. The omitted terms have valuation strictly above $B$ at every admissible tuple and in the weighted Gauss valuation. The same tuple therefore works. Sharpness already occurs in one variable. For a finite packet of nonzero series, including determinants, apply the result to their product to control the \emph{sum} of the specialized valuations by the sum of their fixed-radius Gauss valuations. This conclusion requires a budget on that sum. After a torsion coordinate is substituted the coefficient field may become ramified, so this coordinatewise argument uses the unrestricted exponent.

\section{The exact threshold at arbitrary finite formal height}
\label{sec:height}

The interpolation proof has an elliptic analogue. This matters at a
supersingular prime, where torsion in the formal parameter has a
different distribution from roots of unity. Let $K$ be a finite
unramified extension of $\mathbb Q_p$, let $F$ be a one-dimensional
commutative formal group law over $\mathcal O_K$, and assume its
reduction has finite height $d$. Set $q=p^d$. We use a fixed integral
formal coordinate, work on the open unit disc of $\mathbb C_p$, normalize
$v(p)=1$, and denote multiplication by $p^k$ by $[p^k]$. In this
section $F[p^j]$ denotes its geometric torsion in this disc.

\paragraph{Theorem (finite height).}
Put $s_j=1/((q-1)q^{j-1})$ for $j\ge1$. Fix $h_0\ge1$,
$\rho\ge s_{h_0}$ and $B\ge0$. For every nonzero
$g\in\mathcal O_{\mathbb C_p}[T]$ with $G_\rho(g)\le B$ there is a
point $z$ of exact order $p^h$ with $v(g(z))\le G_\rho(g)$ and
\[
          h_0\le h\le h_0+\lfloor q^{h_0-1}B\rfloor.
\]
The upper exponent is the smallest possible universal choice.
One entire coset of $F[p^{h_0-1}]$ consists of such points of the same
order. These assertions hold also for nonzero power series with integral
coefficients. The coefficients of the tested function can be ramified
even though the formal group has an unramified field of definition.

\paragraph{The torsion metric.}
In $[p](T)$ the coefficient of $T$ is $p$, the coefficients of
degrees $2,\ldots,q-1$ belong to $p\mathcal O_K$, and the coefficient of
$T^q$ is a unit. The last two facts follow from the definition of height
and the unramified base. For $t=v(z)>0$, the unique lowest term has
valuation $qt$ if $t<s_1$, and $1+t$ if $t>s_1$. In the latter
region repeated application increases the valuation by one and cannot
give zero. At $t=s_1$, the image has valuation at least $1+s_1$.
Thus a nonzero point killed by $p$ has valuation exactly $s_1$, and
induction shows a point of exact order $p^j$ has valuation $s_j$.
The first unit coefficient of $[p^j](T)$ occurs in degree $q^j$ by
composition after reduction. Weierstrass preparation gives $q^j$ roots
in the open disc counted with multiplicity. They are simple, since
translation by a torsion point identifies the local zero with the
zero at the origin, where the derivative is $p^j\ne0$.
It follows that there are $(q-1)q^{j-1}$ points of exact order $p^j$.
The series $F(X,[-1](Y))$ is $(X-Y)$ times a series with unit constant
term. Hence $v(x-y)=v(x\mathbin{\ominus_F}y)$ for torsion parameters.
An integral change of formal coordinate with unit linear coefficient
preserves these distances. For the elliptic facts about height see
Boxall and Grant, \emph{Singular torsion points on elliptic curves},
\S4, Lemma 4.1.

\paragraph{Interpolation and the matching obstruction.}
Write $k=h_0-1$, $L=H-k$, and $M=q^k$, with $H\ge h_0$.
The map $[p^k]:F[p^H]\longrightarrow U=F[p^L]$ is onto by the kernel
count above. The allowed points form the inverse image of
$U\setminus\{0\}$. Choose one lift $x_u$ for each nonzero $u\in U$.
A lift of a point of order $p^j$ has order $p^{j+k}$. For different
nonzero $u,w$ this gives
\[
 v(x_u)=v(u)/M,\qquad v(x_u-x_w)=v(u\mathbin{\ominus_F}w)/M.
\]
Each shell in $U$ has total valuation one. Translation permutes $U$,
so for $P(T)=\prod_{u\ne0}(T-x_u)$ we obtain
\[
 G_\rho(P)=L/M,\qquad v(P'(x_u))=L/M-v(x_u).
\]
Indeed the nonzero differences from $u$, after omitting the index
$0$, range over all nonzero group elements except $u$. Since every
$v(x_u)\le\rho$, every Lagrange basis polynomial
$P(T)/((T-x_u)P'(x_u))$ has Gauss valuation zero.

Divide $g$ by $P$. The quotient is integral because $P$ is monic and
integral. If $G_\rho(g)<L/M$, the remainder has the same Gauss
valuation as $g$. Lagrange interpolation then gives
$\min_{u\ne0}v(g(x_u))\le G_\rho(g)$. If every fibre contained a point
where the inequality failed, choosing those points as lifts would
contradict this conclusion. For a bounded integral series, truncate at
degree $D$ with $(D+1)s_H>G_\rho(g)$ and $(D+1)\rho>G_\rho(g)$.
The tail changes neither the Gauss valuation nor the comparison at
any of the finitely many tested points.

Conversely use $g=P$. At a different lift $z$ of the same $u$,
all factors except $T-x_u$ retain their valuations at $x_u$.
For $k>0$, the difference $z\mathbin{\ominus_F}x_u$ is a nonzero element
of $F[p^k]$ and has valuation at least $s_k>s_{k+1}$, larger than
all the inter-fibre valuations. Therefore
$v(P(z))=L/M-v(x_u)+v(z-x_u)>L/M$. At $x_u$ the value is zero.
For $k=0$ every tested point is a root already.
Consequently a full interval of $L$ shells can be blocked at cost
exactly $L/M$. The first integer $L$ with $L/M>B$ gives the asserted
floor and sharp endpoint. For several independent formal coordinates of
the same height the coordinatewise Gauss-minimizer argument gives the
same bound on each exponent; the one-variable example proves sharpness.

\paragraph{A supersingular example.}
The smooth reduction of $E:y^2=x^3-x$ at $3$ has four
$\mathbb F_3$-points. Its trace is zero and its formal group has height
two. With $h_0=2$, $B=1$, and $\rho\ge1/72$, one has $q=9$ and
$H_{\rm opt}=11$. A whole coset of nine torsion points works, with
order between $3^2$ and $3^{11}$. At $H=10$ the extremal polynomial
uses one lift for each of the $9^9-1$ nonzero points of $F[3^9]$.
The comparison concerns functions integral in the local parameter at the
identity. Applying it to a global rational function requires checking that
its expansion has the stipulated integrality and Gauss bound.

\paragraph{Exact tradeoff between degree and valuation budget.}
There is a second sharp parameter that the same interpolation identity
settles at once. Impose $\deg g\le D_0$ as well as $G_\rho(g)\le B$.
A window of $L$ shells works for every such polynomial precisely when
\[
 L>q^{h_0-1}B\qquad\text{or}\qquad q^L-1>D_0 .
\]
Indeed, below the node count the interpolated remainder is $g$ itself,
whatever its Gauss valuation is. When both inequalities fail, $P$
meets both constraints and blocks every point. Equivalently the optimal
upper exponent is
\[
 h_0+\min\bigl(\lfloor q^{h_0-1}B\rfloor,
                    \lfloor\log_q(D_0+1)\rfloor\bigr).
\]
Thus the same polynomial attains the minimal degree and minimal budget of
an obstruction simultaneously. For instance at $p=3$, height one,
$h_0=2$, the degree-26 obstruction in the first theorem is of minimum
degree even if its budget constraint is removed. Degree at most 25
guarantees a good point through exponent 4 at every budget.

\subsection{A valuation cutoff over ramified base fields}
Uniform shell radii were used to express the answer in terms of $d$
and $h_0$. There is an exact replacement when the base $K/\mathbb Q_p$
is ramified. Keep finite reduction height and take $G=F[p^H]$.
For a fixed positive $\rho$, let
\[
 N=\{z\in G:v(z)>\rho\},\qquad
 A_*=\frac{H-\sum_{z\in N,\ z\ne0}v(z)}{|N|}.
\]
Assume $G\setminus N$ is nonempty. Among integral polynomials
that satisfy $v(g(z))>G_\rho(g)$ for every $z\in G\setminus N$, the
minimum possible $G_\rho(g)$ is $A_*$, and the minimum possible degree
is $[G:N]-1$. One polynomial attains both. Below either threshold
there is an entire good coset of $N$.

\emph{Proof.} The formal group law and inverse preserve the condition
$v(z)>\rho$, so $N$ is a subgroup. On each nonidentity coset the
valuation is constant. Adding an element of $N$ changes a parameter by
an amount of valuation greater than $\rho$. For the distinguished
Weierstrass polynomial of $[p^H]$, the product of its nonzero roots
has valuation $H$, by the derivative $p^H$ at zero and the
unit factor of preparation. The total valuation of one representative
from each nonidentity coset is therefore $A_*$. The valuation of
differences between two such representatives depends only on their
difference coset. Translation in $G/N$ now gives
$v(P'(x_u))=A_*-v(x_u)$ for the product $P$ of their linear factors.
The interpolation, arbitrary-representative argument and matching
example above apply verbatim.\hfill$\square$

This formula only needs the multiset of torsion valuations, available
from the Newton polygon; coset counts and the sum over $N$ follow by
counting the slopes. In particular, if the comparison radius includes
all nonzero torsion valuations, $N=\{0\}$ and the obstruction budget
is exactly $H$ over every base. It suffices here to take
$\rho\ge\max_{0\ne z\in F[p]}v(z)$, since multiplication by $p$ strictly
increases the valuation of a nonzero point in the disc until it reaches
zero. This also explains why the closed formula for the unrestricted
height-$d$ order window can change over ramified bases: cutting by
valuation need no longer cut precisely between two orders.


\clearpage
\section{Consequences and remaining questions}
\paragraph{An application to Fourier coefficients of measures.}
Let $\mu$ be a bounded $\mathcal O_{\mathbb C_p}$-valued measure on
$\mathbb Z_p$, and let $a_j=\int\binom{x}{j}\,d\mu(x)$ be its
Mahler moments. For a root of unity $\zeta$ of $p$-power order,
the finite-order Fourier coefficient satisfies
\[
       \widehat\mu(\zeta)=\int_{\mathbb Z_p}\zeta^x\,d\mu(x)
            =\sum_{j\ge0}a_j(\zeta-1)^j .
\]
The identity follows by uniform convergence of the binomial expansion
on $\mathbb Z_p$ when $v(\zeta-1)>0$, first on nonnegative integers
and then by continuity. Thus a single observed moment with
$v(a_j)+j\rho\le B$ guarantees an entire coset of characters of
orders in the sharp window, each with Fourier coefficient of valuation
at most $\min_i(v(a_i)+i\rho)$. In this setting order $p^h$ means
exact additive conductor $p^h$ on $\mathbb Z_p$.
Sharpness persists among finite atomic measures supported on nonnegative
integers: each $T^j=((1+T)-1)^j$ is an integral linear combination
of $(1+T)^r$, $0\le r\le j$, the transforms of point masses.
Accordingly the joint degree and budget obstruction above can be
realized with support in $\{0,\ldots,q^L-1\}$ in the cyclotomic
case $q=p$.

\paragraph{Arithmetic twists.}
An arithmetic use requires an integral interpolation function and a
nonzero witness with a valuation bound. Once these data are given,
the theorem supplies a precise finite range of character orders to
search. For instance, if a determinant has a suitable interpolation
series, the selected determinant value is nonzero. A bound on the
conductor norm, or a numerical Selmer rank bound, requires further
global arithmetic input.
\begin{enumerate}
\item For a fixed ramified coefficient field, which polynomials over
that field attain the unrestricted threshold? The Newton polygon
bound in the first section and the congruences between attainable
conjugate roots impose different restrictions.
\item Over a ramified base for the formal group, can one compute the
threshold for a prescribed order interval when the interval is not a
valuation cutoff? The last formula covers the case of a cutoff exactly.
\item For a positive allowed loss at a fixed Gauss radius, determine the
sharp joint dependence on degree, order interval and tolerance.
\end{enumerate}
\paragraph{Further literature.}
Manjul Bhargava's work on
\href{https://doi.org/10.1006/jnth.1998.2220}{generalized factorials and
fixed divisors} supplies a broader setting for finite-set polynomial
interpolation. Kriz and Nordentoft study torsion specialization for
\href{https://arxiv.org/abs/2310.20678}{horizontal $p$-adic $L$-functions}.
Boxall and Grant describe the elliptic formal groups used here in
\href{https://boxall.pages.math.cnrs.fr/pages-web-mathematiques-de-john-boxall/wwwMaths/SingularFinal.pdf}{their
paper on singular torsion points}, \S4. We claim the improvement over
the stated ProveIt bound. Historical priority in the full interpolation
literature remains unestablished.

\paragraph{Verification.} The companion arithmetic checker computes
cyclotomic polynomials, valuations and values in exact Eisenstein
coordinates. A second checker computes Newton polygons of iterated
Lubin--Tate polynomials and verifies the finite group metric, node
budgets and obstructions in small higher-height examples. The general
claims follow from the proofs above.
\end{document}
